\documentclass[a4paper,12pt]{ctexart}
\usepackage{xcolor}
\usepackage{graphicx}
\usepackage[top=1.8cm, bottom=1.8cm, left=1.8cm, right=1.8cm]{geometry}
\usepackage{array}
\usepackage{hyperref}
\hypersetup{colorlinks=true,linkcolor=blue!70!black}
\linespread{1.35}

\newcommand{\ruby}[2]{#1\textsuperscript{#2}}
\newcommand{\xuehai}[2]{%
  \vspace{0.35cm}
  \noindent
  \fcolorbox{blue!50!black}{blue!5}{%
    \begin{minipage}{0.915\textwidth}
    \textbf{\textcolor{blue!60!black}{学海拾遗：}#1} \\[0.1cm]
    {\small #2}
    \end{minipage}%
  }
  \vspace{0.35cm}
}

\title{\textcolor{blue!70!black}{\Huge\textbf{哲学入门}}}
\author{}
\date{}

\begin{document}
\maketitle
\thispagestyle{empty}

\tableofcontents
\newpage

\section{\textcolor{orange!80!black}{第一课\quad 什么是哲学？}}

\subsection*{一个在广场上追问的老人}

公元前5世纪的雅典。一个六十多岁、塌鼻子、凸眼睛的老人每天在广场上拦住路人——将军、诗人、政客、鞋匠——追问同一个问题："你刚才说'这是正义的'——什么是正义？"

苏格拉底:"你说正义就是欠别人东西要还——对吗？"对方:"当然。"苏格拉底:"假设一个朋友把他的一把刀寄存在你这里。过了一段时间他回来取——但他已经疯了。他拿着这把刀可能伤到自己。你还给他——这算不算正义？"对方愣住了。他把定义从"还欠别人的东西"改为"把属于别人的东西还回去——除非会伤害他"。苏格拉底又追:"那另外一个人来问你要回他的枪。你知道他拿着这枪会去杀一个无辜的人——你该还吗？"

对方一次次往前递定义、一次次被苏格拉底的反例戳回。他最终说的不再是"正义的定义是X"。他说的是\:\textbf{"我不知道。"}\:苏格拉底笑了。\textbf{"你不知道了——这是你和我之间最大的进展。"}

\subsection*{同一时间，东方另一个人也在问}

和苏格拉底差不多同时代，东方的\textbf{孔子}也在追问。但他问的不是"正义是什么"——他在周游列国，对着君主和他的弟子不停追问的是："什么是仁？"子路来问仁，孔子回答:"恭、敬、忠"。冉雍来问仁，孔子说:"出门如见贵宾，使民如承办重大祭祀。"宰我来问仁，孔子直接告诉他"仁很难，做起来很费力"。同一个问题，三种答案——因为孔子在对着三个不同性格的人在不同的生命节点回答。\textbf{苏格拉底相信理性追问可以剖出定义。孔子相信——定义不重要，做出来才是真的。}

\xuehai{哲学、宗教和科学的三角关系}{
  人类回答"世界是什么样"这个问题时，用了三套不同的工具。\textbf{宗教}用启示和经文告诉你世界是神创造的；\textbf{科学}用实验和数学来检验假设（"我猜重的东西掉得更快——让我从塔上扔两个球看看"）；\textbf{哲学}不靠神的启示也不靠实验室——它用逻辑推理和概念分析来追问那些科学还无法做实验的问题（"什么是公正？""自由意志存在吗？"）。

  这三者不是"谁对谁错"的竞争关系。很多伟大的科学家同时是非常虔诚的信徒（牛顿晚年写了一百多万字的圣经研究）。三者的分工大致可以这样理解：科学管"事实是什么样的"（fact），哲学管"我们应该怎么理解这些事实"（interpretation），宗教管"这一切的意义是什么"（meaning）。
}

\subsection*{哲学不负责给你标准答案}

数学考试问7+5等于多少——12。错的人是有计算错误。但哲学的核心问题没有最终统考答案。\textbf{哲学只是帮你把自己脑子里那套隐藏的判断标准挖上来——擦干净——摆在你面前。}苏格拉底死前的一句名言:"未经审视的生活是不值得过的。"

\subsection*{想一想}

\begin{enumerate}
  \item 苏格拉底一直追问到对方说"我不知道"。为什么他反而认为这是进展？知道"自己不知道"比你随口扔出的一个定义更诚实——也更危险吗？
  \item 孔子对三个学生给出了完全不一样的"仁"的回答。他在偷懒吗？还是"因人而异"本身就比"给所有人同一个正确公式"更难、更接近真实？
  \item 宗教用经文、科学用实验、哲学用追问。你今天遇到"为什么活着"这个问题时——你更习惯于调用哪一种工具？
\end{enumerate}

\newpage
\section{\textcolor{orange!80!black}{第二课\quad 我怎么知道我知道？}}

\subsection*{你确定你面前有一张桌子？}

你现在把手放在一张硬硬的平坦表面上。你非常确信——这是一张桌子。但你为什么如此确信？\textit{因为你看见了——它看上去是棕色的，有一条条木纹似的东西。}你的眼睛输入了颜色和形状的信号。\textit{因为你摸得到——手掌得到坚实的阻力。}你的触神经把压力数据传进了大脑。你并没有直接摸到"桌子"本身。你接收到的只是一串神经信号——颜色、阻力、纹理——大脑把这些信号拼成了一个叫"木桌"的图案。

17世纪的法国哲学家笛卡尔做过一个更极端的锻炼：\textbf{我能不能怀疑一切？}我眼前这个"窗外有路灯"可能只是我做的一个极逼真的梦。我的全部感官都可能被一个看不见的恶魔幻觉操纵。但有一件东西笛卡尔认为自己可以确信：\textbf{如果我能怀疑一切——那至少"正在怀疑的我自己"是存在的。}即使感官全作假——"思考"这个动作仍然进行着。\textbf{我思，故我在。}

\subsection*{教室里的认知实验}

你们班今天做一个实验。黑板上写了五个五位数。老师让全班默记一分钟。然后她请四个同学去走廊等着，剩下的班级把黑板上的第三和第五个数调换了位置。把四个同学叫回来——问他们"这些数字跟刚才是一样的吗？"多数人没发现。\textbf{你的记忆不是一台录像机。它是一台不断的在编辑和重写旧片段的大脑——而你感觉不到每一刀的修改。}

\subsection*{手机上的一段视频——你一定信吗？}

你刷到一个短视频：某位知名演员对着镜头说："我宣布——今晚开始我会把所有粉丝的留言一条一条亲自回复！"声音就是他，嘴型完全对得上。几秒钟之后字幕弹出：\textbf{"本视频由AI合成。"}\textbf{如果你的眼睛和耳朵可以被一个手机上的软件轻松欺骗——"亲眼所见"还管用吗？}

\xuehai{可证伪性——科学和迷信之间唯一的隔离墙}{
  科学哲学家卡尔·波普尔在20世纪提出了一个简洁的判断标准：\textbf{一个理论要算是科学的，它必须能被一个可能的观察推翻。不能在任何可想象的情况下被证明为假的理论——不管听起来多合理——都不属于科学。}

  弗洛伊德说"所有小男孩都潜意识里爱恋母亲并嫉妒父亲"。如果一个小男孩表现出恋母——证实了弗洛伊德。如果他完全不表现出恋母——弗洛伊德说"他的恋母情节被压抑了"。无论如何都"对"——这就是不可证伪。不是说它一定是错的——是说科学没法检验它。区分科学和伪科学不在于"对错"，在于\textbf{理论本身有没有留出一扇可以被事实推翻的门}。
}

\subsection*{想一想}

\begin{enumerate}
  \item 教室里的数字被改了，很多人没发现。如果几分钟的短期记忆都能被偷偷编辑——你对更早以前的事情有多少把握？
  \item 你看了一段AI换脸的视频。如果字幕没标——你会信吗？当眼睛和耳朵很可能不再可靠时，你还有什么方法来证明一段影像是不是真的？
  \item "所有失恋的人都会说'我早就料到会分手'。"如果你追问一个失恋的朋友——她确实这样说了。可证伪吗？如果她在热恋时说过"我们绝对不可能分"——哪一个才是她的"真实想法"？
\end{enumerate}

\newpage
\section{\textcolor{orange!80!black}{第三课\quad 我该怎么做？}}

\subsection*{电车难题}

一辆失控的有轨电车正疯狂地冲向分岔道口的绑着的五个人——他们躲不开。你正站在道岔拉杆旁：如果你推下拉杆，列车会拐入另一条侧线——那条侧轨上绑着一个人。推——还是不推？

\subsection*{尺子一：功利主义——算账}

边沁说：\textbf{法律的正确性由它产出的幸福总量大小来衡量。}把每个人的快乐和痛苦都当作平等的计算单位——尽可能产生最大多数的最大幸福。照着这把尺子：死一个而救五个——数字账上是+4，侧线是正确选项。但把人本身当可加计的算术单元——有没有让你不舒服？

\subsection*{尺子二：康德——别把人当手段}

康德强烈反对。他说：\textbf{不可通过把一个无辜的人当作一件工具来达成善的目的。}如果你推拉杆，你做出了一个蓄意行为对准一个无辜的人。在康德的眼里——任何理性的生命都有同等不可被当作非人手段的价值。

\subsection*{尺子三：亚里士多德——问"我要成为怎样的人"}

亚里士多德觉得前面两种吵法都把问题问错了。重要的不是"这条规则对不对"——而是\textbf{我希望自己长成什么样的人？}\:勇气、节制、公正感、合情理——这些德性都是需要反复做、直到变成第二天性的习惯。\textbf{拉不拉杆不是算出来的——是你选择做哪一种人。}

\subsection*{一张真的病床——三把尺子全撞上了}

2020年初，意大利贝加莫。新冠肺炎病人涌进医院的速度远远超过了呼吸机的数量。医生们被迫做了判断：一台呼吸机——两个病人在床边等。谁先用？功利主义在纸上说："选存活概率最高的那个人。"康德在脑中说："你不能把任何一个等在床边的病人当作'次选生命'。"亚里士多德低声问："这个夜晚——你自己想成为哪一种医生？"\textbf{三把尺子没有一个能把另外两把完全压死——但躺在床上的病人的氧气面罩等不了哲学家继续吵。}

\xuehai{道德哲学的三个流派——一句话总结}{
  \textbf{后果论}（功利主义是最著名的代表）：判断一个行为对不对，看它产生的后果。好结果多的就是对的。\textbf{义务论}（康德是代表）：有些行为本身就是错的——不管它可能带来多大好处。不能把任何人当作纯粹的工具。\textbf{德性伦理学}（亚里士多德是源头）：不要老问"我应该遵守哪条规则"，问"我想成为什么样的人"。

  这三种立场在大脑里对应的直觉不同：后果论调用的是你的计算脑区（"五个大于一个"），义务论调用的是你的道德情感（"不能推人——推人不对"），德性论调用的是你的身份认同（"我是不是那种推人的人"）。
}

\subsection*{想一想}

\begin{enumerate}
  \item 医院里没有足够的呼吸机——三把尺子都不能完美回答。你有第四把尺吗？
  \item 你朋友偷偷把班费拿去买了自己的新球鞋。你和他是最好的朋友。告诉老师还是不告诉？康德的绝对命令在友谊面前崩溃吗？
  \item 你更偏向功利主义、康德还是亚里士多德？找一个你最近做过的决定——用这三把尺子各量一次，看看你的决定能不能同时通过三把尺子的检验。
\end{enumerate}

\newpage
\section{\textcolor{orange!80!black}{第四课\quad 自由意志——我的选择是我的选择吗？}}

\subsection*{"我决定按了"——你的脑子比你早}

科学家让被试坐在一台计时器前，随意决定什么时候按一下左键或右键并记录下自己"决定的瞬间"。同时他们用脑电帽套在被试头顶。\textbf{运动皮层的预备激活信号，在受试者主观报告"我刚选择了"之前的0.3到数秒就已经在脑电图上出现了。}你的大脑在你意识到"我做了决定"之前——已经在悄悄开始准备动作了。

\subsection*{你的基因说——你更可能暴力}

有一种基因叫MAOA。一种低活性的变体——有人管它叫"战士基因"——本身不能"让你变暴力"。但它在一个特定条件下会显著放大风险：携带这种基因变体的人如果在童年时期遭受过严重的虐待或忽视，长大后犯暴力罪的概率远高于普通人群。现在把这个人的档案放在法官面前。\textbf{他选择去打劫的那一刻——有多少成分是他自己在选的？}

\subsection*{水壶里的兴奋剂}

你的教练在你比赛前递给你一瓶水。你不知道他在水里放了兴奋剂。你喝了——赛后尿检呈阳性。你被禁赛两年。反兴奋剂体系给了一个答案：\textbf{严格责任}——不管你是不是故意的，只要你的身体里出现了被禁止的物质，你就必须承担责任。

\xuehai{决定论和自由意志——三个阵营}{
  哲学家把这个问题拆成了三个主要站位。\textbf{硬决定论}：世界上每一个事件——包括你的每一个"选择"——早在宇宙大爆炸那一刻就由物理因果链条决定了。你觉得你自由——只是因为你没看到所有的链条。\textbf{自由意志论}：人的意识有某种不受物理定律完全支配的"自由选择"能力——和决定论不能同时为真。\textbf{相容论}（中间派）：决定论和自由意志可以同时为真——只要"自由"被定义为"你的行动是你自己的欲望和信念导致的、没有外部强制"，哪怕你的欲望本身也是被因果链决定的。

  这个问题不只是一个哲学游戏——它直接决定着一个社会的法律该怎么写。如果一个人的犯罪行为是被基因和童年的虐待完全"决定"的——惩罚他还有道德意义吗？还是说惩罚不是为了"道德谴责"，而是为了防止他再犯、警示别人的"实用工具"？
}

\subsection*{想一想}

\begin{enumerate}
  \item 脑电图在你的"意志"出现之前就显示出信号——肚子饿了自动反应和你做的"选择"之间有什么真正不同的地方？
  \item 运动员喝了教练递过来的水而不知道里面有药——该不该全责禁赛？如果"不知道"就可以免罚，会不会每个人都把药藏在另一个人的水壶里？
  \item 相容论说"只要没有外部强制就算自由"。你从小在一个只允许吃素的家庭长大——你从来没自己决定过"要不要吃肉"。你是被"决定"了——还是"自由"地选择了素食？
\end{enumerate}

\newpage
\section{\textcolor{orange!80!black}{第五课\quad 什么是好的生活？}}

\subsection*{体验机器——你最快乐的选择}

哈佛哲学家罗伯特·诺齐克在1970年代提出了一个思想实验。你觉得你能不能拒绝进入一个\textit{体验机器}——它是一个巨大的水缸把你全身浸泡，电极连接你大脑的全部愉快回路。你一进入——从你这一面体验到的是连续不断的极致幸福和快乐。\textbf{诺齐克问所有人的那一个问题是："你要进吗？"}绝大多数的被试回答——不进。为什么？因为\textbf{人不只想要"感到快乐"}。我们还想让快乐的来源自己是真实的。\textbf{意义不可以被软件代劳。}

\subsection*{三条河，同一个问题}

\textbf{西方·亚里士多德：}幸福不是"感到开心"。它是一条生命按照德性进行的活动——活动是动词。\textbf{中国·孔子：}人生的坐标不是你自己。心安——不是在于你赚了多少钱，而是在你对得起所有你该珍视的人。\textbf{印度·释迦牟尼（佛陀）：}痛苦的来源不是命运。是执著。抓着不放——苦从手指缝渗进去。松开手，不等著，路空了。

\subsection*{考上了所有人都想去的学校——然后呢？}

你认识一个学姐。初三那年全年每一晚几乎没在凌晨一点前睡过。七月份她收到了那所所有人都在盯着看的学校的录取短信。第四天她把自己锁在房间里。"所有力气全用在这一场试上了——现在考试考完了——\textbf{我不知道接下来该为谁拼命。}"

\textbf{诺齐克会说——这恰恰证明了你不想只活在"体验机器"里。你需要的不是一串快乐的信号——而是下一件值得你亲手去做的、困难的、能让你感觉今天的自己和昨天有区别的事。}

\xuehai{幸福感的两条路径——享乐型和意义型}{
  心理学家把幸福感拆成了两种：\textbf{享乐型幸福}追求的是愉悦感和满足感——美食、好天气、好看的电影。\textbf{意义型幸福}来自于感到自己的生命在和某种比自己更大的事物相连——养大一株花、帮到一个朋友、做出一件自己满意的东西。

  这两种幸福在某一个点之前是互补的——但过了某个拐点之后，追加享乐不会带来相应的幸福增量。多买一件玩具——当天快乐——第二天玩具躺在床底下积灰。但你把那个周末花在了帮同学弄懂他不会的一道题上——那一小片"被需要"可以在你的记忆里留几个月。你不需要在两者之间二选一——但你需要注意你每天的"幸福预算"里——享乐花了多少、意义存了多少。
}

\subsection*{想一想}

\begin{enumerate}
  \item 如果你妈可以把你的意识一键上传到"永远快乐永远不会被批评"的学校——你愿不愿意注销原来的世界？
  \item 孔子、亚里士多德和佛陀指向了三个截然不同的方向。你今天的日常——更偏向哪一条河？
  \item 你最近一次感到"这件事真难但我做完后觉得特别值"是什么时候？和"吃了一个好吃的冰淇淋"相比——哪一种快乐在你身体里留得更久？
\end{enumerate}

\newpage
\section{\textcolor{orange!80!black}{第六课\quad 什么是公正？}}

\subsection*{无知之幕}

美国哲学家约翰·罗尔斯设计了一个著名的思想实验：假设你站在自己出生之前。\textbf{你不知道自己会降生在富裕家庭还是贫穷家庭，不知道自己是男生还是女生，不知道自己的头脑是自动数理型还是音乐耳朵。}在这层"无知之幕"后面——和所有将降生的人一起设计社会的基本规则。你会写一条——"把所有好东西全归给最聪明的那个人"？你不会——因为"最聪明的那个人"可能不是你。\textbf{社会和经济的不平等——只有在它们能改善最底层处境时才是被允许的。}

\subsection*{诺齐克的回怼——不要强迫施舍}

罗伯特·诺齐克——就是那个设计了体验机器的同一个人——对罗尔斯的再分配逻辑发了严厉的批评：如果篮球巨星张伯伦每晚打比赛的门票全是观众自愿掏25美分单独付给他的——他合法赚到了一座小山式的钱。国家现在可以因为你罗尔斯说的"最弱者优先"把他合法挣来的钱强行分给陌生人吗？诺齐克说：不可以。

\subsection*{平板电脑该发给谁？}

学校新到了一批平板电脑——数量只够发一个班。选项A：发给尖子班——能最大程度利用平板的功能。选项B：发给学习最困难的班——他们每天借不到家里的笔记本电脑完成网上作业。你坐在那间会议室里，\textbf{你要怎么投你那一票？}

\xuehai{税收和再分配——三把不同的"公平秤"}{
  几乎所有政治哲学关于"公正"的争论最终都会掉进同一个问题里：\textbf{政府应该从收入更高的人那里拿走多少钱、花在收入更低的人身上？}

  \textbf{边沁/密尔的功利主义}：能增加全社会总幸福就是合理的——如果从富人那里拿走一块钱给他带来的痛苦小于这一块钱给穷人的快乐——拿走。\textbf{罗尔斯的差异原则}：不平等只有在它让最底层的人也获益时才被允许——给外科医生高薪让更多人被治好比平均分配但要不到半小时手术更公平。\textbf{诺齐克的"最小国家"}：国家唯一的任务是保护生命和财产、执行合同——做了这三件之外的事都是强迫。最穷的人得到的最"公正"的帮助不是国家的再分配——是富人的自愿捐赠。

  三把秤称同一斤米，三个刻度。\textbf{没有一个能让另外两个彻底闭嘴——但每一次预算讨论和每一次选战，你都看见这三把秤的指针在公众辩论中来回撕扯。}
}

\subsection*{想一想}

\begin{enumerate}
  \item 平板分给了尖子班——困难班的同学连作业都交不上。困难班的人有没有自己被"合法的放弃"？
  \item 反过来平板分给困难班——尖子班的学生说"我们本可以做出更好的项目"。另一个方向的怨恨——公平付给哪一边的损失？
  \item 你自己选了出生的家庭吗？你的天赋和家境有一部分不是你自己挣的——它们是你"应得"的还是你"抽到的"？如果天赋不是自己挣的——高天赋的人该不该被多征税？
\end{enumerate}

\newpage
\section{\textcolor{orange!80!black}{第七课\quad 逻辑与谬误——如何不被忽悠}}

\subsection*{稻草人、假两难、滑坡}

"我觉得这周的语文作业可以少一点点——因为背诵三篇长课文需要时间。""\textbf{你就是说学语文没用吧！}"——稻草人。"你不买这个保险——你的以后一辈子的医药费谁付？"——假两难：把中间选项全删光。"今天早上迟到五分钟老师没罚——下一步全班每天迟到——再下一步全校迟到——再下一步全国的学校全乱套了！！"——滑坡：缺了每一个台阶之间的证明。

\subsection*{相关不等于因果——冰淇淋和溺水的那个笑话}

某市警察署发布数据：冰淇淋销量上涨的月份溺水死亡人数同时上涨。有人跑出来开记者会——"冰淇淋是一项危险的杀人产品！"同涨同落，没错。但原因是什么？是夏天。热——更多人买冰淇淋——同时更多人去游泳。\textbf{数据常常告诉你A和B一起动了——剩下的你必须用自己的脑子去剥。}

\subsection*{动手造一个"好看的片子"}

请你自己试着写一小段假新闻——只用三四句话——\textbf{用上任意一种你今天学到的谬误手法}——但要写得让别人第一眼读过去会真的停下来信一秒。

"研究表明：每天戴耳机听歌超过三小时的学生期中成绩平均低11分。耳机正在毁掉这一代中学生的学业。"\textit{捏这条消息的人在你的课后习题边打了勾：相关不等于因果、假两难、没有控制变量。}

\xuehai{批判性思维的三个基本问题}{
  不论你面对的是新闻标题、短视频解说还是社交帖子，有三把钥匙可以捅进任何一段信息的锁孔里：

  \textbf{第一，谁在说、为谁的利益说？}一条"每天喝八杯水有益健康"的建议——发布者是自来水公司、矿泉水品牌、还是独立的医学研究机构？信源的利益立场不影响这句话的真假——但你需要在脑内举一个牌子提醒自己"这一方有一台自己的秤"。

  \textbf{第二，证据链完整吗？}冰淇淋vs溺水——有相关，没有因果证据。"自从换了校长，一本率涨了8\%——新校长真厉害"——同期还翻修了空调、换了新版教材、把晚自习延长了四十分钟。你被给了一个"因为"——同时给你隐藏了四条其他的"因为"。

  \textbf{第三，你的情绪在这一秒被用来做什么？}"不转不是中国人""是家长就转"——这是把你的恐惧和愤怒当作一个按钮来按。\textbf{当你感到自己的愤怒或恐慌忽然被一句话同时点燃时——先停下来，数三秒，问：这句是要让我想，还是要让我转？}
}

\subsection*{想一想}

\begin{enumerate}
  \item 找出你最近在群聊或短视频里看到的一句听起来非常有道理、但仔细一拆有一个明显谬误的话。把它写下来，标注它犯了哪一类错误。
  \item "自从新校长上任，学校的一本率上涨了8\%——新校长就是比旧校长强。"同时期这所学校还翻修了所有教室的空调、换了新版的教材、并且把晚自习延长了四十分钟。哪一个变量在提分——是你确信的"那个"吗？
  \item 你在社交媒体上看到一条消息让你愤怒得几乎要转发——你停了三秒。请自问：这条消息给我的是关键证据还是情绪按钮？信源是谁？有没有隐藏的"因为"？
\end{enumerate}

\vspace{0.5cm}
\begin{center}
\textcolor{blue!70!black}{\large —— 七课之后，你不会拿到一份叫"哲学入门"的标准答案。但你手里多了一副能切开谎话和推自己一把看清自己的工具。——}
\end{center}

\end{document}
